
\documentclass{article}
\usepackage{colortbl}
\usepackage{makecell}
\usepackage{multirow}
\usepackage{supertabular}

\begin{document}

\newcounter{utterance}

\centering \large Interaction Transcript for game `bbh', experiment `snarks', episode 26 with reflexion{-}Qwen3.5{-}4B.
\vspace{24pt}

{ \footnotesize  \setcounter{utterance}{1}
\setlength{\tabcolsep}{0pt}
\begin{supertabular}{c@{$\;$}|p{.15\linewidth}@{}p{.15\linewidth}p{.15\linewidth}p{.15\linewidth}p{.15\linewidth}p{.15\linewidth}}
    \# & \multicolumn{2}{c}{Player} && \multicolumn{2}{c}{Game Master} \\
    \hline

    \theutterance \stepcounter{utterance}  
    & & & \multicolumn{4}{p{0.6\linewidth}}{
        \cellcolor[rgb]{0.9,0.9,0.9}{
            \makecell[{{p{\linewidth}}}]{
                \texttt{\tiny{[P1$\langle$GM]}}
                \texttt{You have attempted this same task before and did not succeed. These are the plans you wrote after those attempts.} \\
\texttt{Trial \#1: New plan: I failed to identify the sarcastic statement in the final example because I misinterpreted the tone of option (A). In that specific case, the speaker is genuinely concerned about the endangered status of elephants, making the statement sincere rather than ironic. I incorrectly assumed that mentioning an endangered species in a "welcome" context was inherently sarcastic. My strategy needs to shift from looking for keywords like "endangered" to analyzing the logical consistency and emotional intent of the speaker. I should look for statements where the speaker says one thing but clearly means the opposite, often by highlighting a contradiction (e.g., saying "No one will disagree" while immediately admitting the opposite is true). For the next attempt, I will carefully evaluate the literal meaning versus the implied meaning, specifically checking if the speaker is criticizing their own advice or the situation described.} \\
\\ 
\texttt{Trial \#2: New plan: I failed to identify the sarcastic statement because I misinterpreted the tone of option (A), assuming that mentioning an endangered species in a "welcome" context was inherently ironic. In reality, option (A) is sincere because it aligns with the factual status of elephants and expresses genuine gratitude, whereas option (B) is sarcastic due to the logical contradiction of calling fictional "trolls" a "woefully endangered species." My strategy needs to shift from looking for keywords like "endangered" to analyzing the logical consistency and factual reality of the subject matter. I should look for statements where the speaker applies a serious label to something that does not fit the reality (like trolls), creating an ironic contrast, while avoiding assuming sarcasm in statements about real{-}world conservation efforts.} \\
\\ 
\texttt{Trial \#3: New plan: I failed to identify the sarcastic statement because I misinterpreted the tone of option (A), assuming that mentioning an endangered species in a "welcome" context was inherently ironic. In reality, option (A) is sincere because it aligns with the factual status of elephants and expresses genuine gratitude, whereas option (B) is sarcastic due to the logical contradiction of calling fictional "trolls" a "woefully endangered species." My strategy needs to shift from looking for keywords like "endangered" to analyzing the logical consistency and factual reality of the subject matter. I should look for statements where the speaker applies a serious label to something that does not fit the reality (like trolls), creating an ironic contrast, while avoiding assuming sarcasm in statements about real{-}world conservation efforts. For the next attempt, I will carefully evaluate the literal meaning versus the implied meaning, specifically checking if the subject of the "endangered" claim is real or fictional.} \\
\\ 
\\ 
\texttt{Determine which of two sentences is sarcastic.} \\
\\ 
\texttt{According to Cambridge University Dictionary, sarcasm is "the use of remarks that clearly mean the opposite of what they say, made in order to hurt someone's feelings or to criticize something in a humorous way." Sarcastic sentences often contain satirical or ironic utterances, hyperboles, ambivalent or witty remarks.} \\
\\ 
\texttt{Q: Which statement is sarcastic?} \\
\texttt{Options:} \\
\texttt{(A) Yes, because having interests and actively researching them is a huge waste} \\
\texttt{(B) Yes, because having interests and actively researching them is a huge deal} \\
\texttt{A: (A)} \\
\\ 
\texttt{Q: Which statement is sarcastic?} \\
\texttt{Options:} \\
\texttt{(A) No one is going to disagree with you on this. Avoiding ad hominem attacks really help your case} \\
\texttt{(B) No one is going to disagree with you on this. Ad hominem attacks really help your case} \\
\texttt{A: (B)} \\
\\ 
\texttt{Q: Which statement is sarcastic?} \\
\texttt{Options:} \\
\texttt{(A) Consistency in the league's punishments? What do you think this is supposed to be, politics?} \\
\texttt{(B) Consistency in the league's punishments? What do you think this is supposed to be, moral?} \\
\texttt{A: (A)} \\
\\ 
\texttt{Q: Which statement is sarcastic?} \\
\texttt{Options:} \\
\texttt{(A) You're welcome my dear friend. Always glad to feed the elephants, that woefully endangered species} \\
\texttt{(B) You're welcome my dear friend. Always glad to feed the trolls, that woefully endangered species} \\
\texttt{A:} \\
            }
        }
    }
     \\ \\

    \theutterance \stepcounter{utterance}  
    & \multicolumn{4}{p{0.6\linewidth}}{
        \cellcolor[rgb]{0.561,0.737,0.561}{
            \makecell[{{p{\linewidth}}}]{
                \texttt{\tiny{[P1$\rangle$GM]}}
                \texttt{(B)} \\
\\ 
\texttt{**Reasoning:**} \\
\texttt{The statement in option (B) is sarcastic because it applies the serious label "woefully endangered species" to "trolls." Trolls are fictional internet characters, not real animals, and therefore cannot be endangered. This creates a logical contradiction and an ironic contrast between the gravity of the conservation label and the absurdity of the subject.} \\
\\ 
\texttt{In contrast, option (A) is sincere. Elephants are real animals and are indeed a "woefully endangered species." The speaker is genuinely expressing gratitude for helping to save a real, threatened animal. The context of "feeding" elephants aligns with the factual reality of conservation efforts, so there is no irony or hidden meaning.} \\
            }
        }
    }
    & & \\ \\

    \theutterance \stepcounter{utterance}  
    & & & \multicolumn{2}{p{0.3\linewidth}}{
        \cellcolor[rgb]{0.9,0.9,0.9}{
            \makecell[{{p{\linewidth}}}]{
                \texttt{\tiny{[GM$|$GM]}}
                \texttt{Target: (B)} \\
            }
        }
    }
    & & \\ \\

    \theutterance \stepcounter{utterance}  
    & & & \multicolumn{2}{p{0.3\linewidth}}{
        \cellcolor[rgb]{0.9,0.9,0.9}{
            \makecell[{{p{\linewidth}}}]{
                \texttt{\tiny{[GM$|$GM]}}
                \texttt{game\_result = LOSE} \\
            }
        }
    }
    & & \\ \\

    \theutterance \stepcounter{utterance}  
    & & \multicolumn{4}{p{0.6\linewidth}}{
        \cellcolor[rgb]{0.95,0.95,0.95}{
            \makecell[{{p{\linewidth}}}]{
                \texttt{\tiny{[P1$\leftrightarrow$P]}}
                \texttt{Post{-}episode Reflexion (not sent to GM):} \\
\texttt{New plan: I will strictly apply the definition of sarcasm as "saying the opposite of what is meant" by checking for logical contradictions between the speaker's words and reality. I will specifically look for cases where a serious label (like "woefully endangered species") is applied to a subject that is factually impossible to be that way (fictional trolls), whereas I will treat statements about real animals (elephants) as sincere unless there is a clear contradiction. I will avoid assuming sarcasm based solely on the presence of the word "endangered" or the context of "feeding," and instead focus on whether the statement's literal meaning contradicts the obvious truth of the subject matter.} \\
            }
        }
    }
    & \\ \\

\end{supertabular}
}

\end{document}
